% Kapitel 5/8 — Masse-Konstitution (EABC / Hurwitz)
\begin{axiom}[Masse nur zwischen sichtbaren EABC-Zwischenstrukturen]
Sei $H$ die Hurwitz-Ordnung mit $E,A,B,C\leftrightarrow 1,i,j,k$ und
$\Pi_\Gamma=\frac14(I+\Gamma+\Gamma^2+\Gamma^3)$ der Schalen-Kollaps aus dem holonomen
Rundweg $\Gamma=R\cdot K$.
Eine Massengröße $m$ ist nur auf einem Zwischenobjekt
$z\in\mathcal{Z}$ (aktive Kante, Schwung $u=\overline{q_k}q_{k+1}$, neutraler Vierling
oder Bindungsdipol $B\circ C=\pm A$) definiert, wenn für das zugehörige Aggregat $q\in H$:
\begin{enumerate}
  \item $\Pi_\Gamma(q)\neq 0$ \emph{(elektrodynamisch sichtbar im glatten Schalenanteil)};
  \item der tragende Übergangsgraph erlaubt Zweiweg-Kopplung
        (bidirektional wie $T+T^{-1}$ bzw.\ reversibler Schwung);
  \item $N(\Pi_\Gamma(q))\ge 1$ in der Hurwitz-Norm.
\end{enumerate}
Fluktuationen in $V_{\lambda=0}=\ker(\Pi_\Gamma)$ tragen keine Masse;
arithmetische Resonanz $q\in V_{\lambda=0}\cap H\pi_p$ ist von der Masse-Konstitution zu trennen.
\end{axiom}
